\documentclass[11pt]{article}
\usepackage[T1]{fontenc}
\usepackage{lmodern}
\usepackage[margin=0.92in]{geometry}
\usepackage{amsmath,amssymb,amsthm,mathtools}
\usepackage{booktabs,longtable,array}
\usepackage{xcolor,microtype,fancyhdr,enumitem}
\usepackage[hyphens]{url}
\usepackage{xurl}
\usepackage{listings}
\usepackage[hypertexnames=false,colorlinks=true,linkcolor=blue!45!black,
 citecolor=blue!45!black,urlcolor=blue!45!black]{hyperref}
\setlength{\headheight}{14pt}
\setlength{\emergencystretch}{3em}
\setlist{itemsep=3pt,topsep=5pt}
\lstset{basicstyle=\ttfamily\footnotesize,breaklines=true,
 columns=fullflexible,keepspaces=true,frame=single,rulecolor=\color{black!18}}
\pagestyle{fancy}\fancyhf{}
\fancyhead[L]{\small Constructive Lagrange inversion}
\fancyhead[R]{\small iPRA and finite coefficient algebra}
\fancyfoot[C]{\thepage}

\newtheorem{theorem}{Theorem}[section]
\newtheorem{lemma}[theorem]{Lemma}
\newtheorem{proposition}[theorem]{Proposition}
\newtheorem{corollary}[theorem]{Corollary}
\newtheorem{claim}[theorem]{Claim}
\theoremstyle{definition}
\newtheorem{definition}[theorem]{Definition}
\newtheorem{remark}[theorem]{Remark}
\newtheorem{audit}[theorem]{Audit finding}

\newcommand{\N}{\mathbb N}
\newcommand{\Z}{\mathbb Z}
\newcommand{\Q}{\mathbb Q}
\newcommand{\C}{\mathbb C}
\newcommand{\ipra}{\mathsf{iPRA}}
\newcommand{\PR}{\mathsf{PR}}
\newcommand{\fps}[1]{#1[\![X]\!]}
\newcommand{\Res}{\operatorname{Res}}
\newcommand{\coeff}[2]{[X^{#1}]\,#2}
\newcommand{\ord}{\operatorname{ord}}
\newcommand{\abs}[1]{\left|#1\right|}
\newcommand{\norm}[1]{\left\lVert#1\right\rVert}
\newcommand{\code}[1]{\texttt{#1}}
\newcommand{\A}{\mathcal A}
\newcommand{\B}{\mathcal B}
\newcommand{\T}{\mathcal T}

\title{Lagrange inversion in intuitionistic primitive recursive arithmetic\\
Finite residue calculus, coefficient formulae, analytic lifting, and applications}
\author{A constructive proof package prepared for Vladimir Reshetnikov}
\date{November 15, 2026}
\hypersetup{pdftitle={Lagrange inversion in intuitionistic primitive recursive arithmetic},
 pdfauthor={Prepared for Vladimir Reshetnikov},
 pdfsubject={Constructive proofs, exact coefficient checks, and a Lean formalization roadmap}}

\begin{document}
\maketitle

\begin{abstract}
We give a constructive development of the mathematical assertions on the
Wikipedia page ``Lagrange inversion theorem.''  The algebraic core is stronger
than a merely analytic argument: for every requested truncation order it is a
finite identity in the ring of truncated series.  The unique solution of
\(G=Z\Phi(G)\) is produced by primitive iteration, and a finite Laurent-series
change-of-variables lemma proves both standard Lagrange--B\"urmann coefficient
formulae.  All sums, products, compositions, derivatives, determinants, Bell
polynomials, and coefficient extractions are bounded algorithms.  Consequently
their correctness is formalizable in first-order intuitionistic primitive
recursive arithmetic using quantifier-free induction.

From the core theorem we derive the inverse-series formula around an arbitrary
base point, the Bell-polynomial expression, Good's multivariate formula, and a
ramified/Puiseux version when the first derivative vanishes.  We then prove the
page's algebraic equation example, the two displayed Lambert \(W\) expansions,
the Catalan-number binary-tree count, and the coefficient transformation used
in Laplace--Erd\'elyi asymptotics.  A quantitative contraction argument supplies
an explicitly positive analytic inverse radius for primitive-recursively named
analytic data.  Exact radii are established for the algebraic and Lambert
examples.

The source audit also detects a normalization mismatch in the page's
associahedron display.  We prove the standard ordinary-coefficient face formula,
translate it correctly to the exponential convention used immediately before
it, and give a degree-two check showing that the displayed formula cannot be
right as written.

The archive contains this paper, a detailed Lean blueprint tied to Mathlib
v4.32.0 and ProveIt's \code{ConciseIPRA} interfaces, and an exact-rational
executable verifier.  A completed kernel-checked Lean formalization is not
claimed.
\end{abstract}

\newpage
\begingroup\small\tableofcontents\endgroup
\newpage

\section{Scope, interpretation, and result audit}\label{sec:scope}

\subsection{What it means to prove an analytic statement in \(\ipra\)}

First-order \(\ipra\) quantifies over natural numbers.  It has no object
variables ranging over arbitrary complex functions, completed infinite
sequences, or open subsets of \(\C\).  Therefore an assertion such as ``every
analytic germ with nonzero derivative has an analytic inverse'' must be read as
a theorem schema for effectively supplied data.

A complex number is represented by a regular primitive-recursive sequence of
Gaussian rationals.  An analytic germ is represented by a fixed
primitive-recursive coefficient procedure, or by a fixed primitive-recursive
rational evaluation procedure, together with explicit rational bounds and
moduli.  Every occurrence of convergence has a supplied primitive-recursive
rate.  Once these procedures are substituted, the theorem is an ordinary
first-order arithmetic sentence.  This is the same constructive convention
used by the existing ProveIt Taylor/Cauchy packet \cite{proveitTaylor}.

The formal-power-series core is even simpler.  For a requested coefficient
bound \(N\), a series is the finite vector
\[
 (a_0,a_1,\ldots,a_N),
\]
and every operation is performed modulo \(X^{N+1}\).  A uniform theorem in
\(N\) is a theorem about natural-number codes of finite vectors and bounded
folds.  No completed infinite object is needed inside the object theory.

\begin{remark}[No hidden universal evaluator]
A theorem can be uniform in a numerical parameter and in the values returned
by a fixed primitive-recursive coefficient procedure.  It cannot quantify over
``all primitive-recursive programs'' via a primitive-recursive universal
evaluator: the usual diagonal function would contradict universality.  Thus the
function symbols and their correctness hypotheses are fixed before the
first-order theorem is formed.
\end{remark}

\subsection{Audit of the page's mathematical assertions}

Table~\ref{tab:audit} records what is proved and in which sense.  ``Finite''
means a coefficient identity in a truncated ring, uniformly primitive
recursive in the truncation order.  ``Named analytic'' means the corresponding
statement for primitive-recursively named complex data with explicit bounds and
moduli.

\begin{longtable}{>{\raggedright\arraybackslash}p{0.28\textwidth}
                  >{\raggedright\arraybackslash}p{0.17\textwidth}
                  >{\raggedright\arraybackslash}p{0.47\textwidth}}
\caption{Status of the mathematical statements on the reference page.}
\label{tab:audit}\\
\toprule
Statement & Status & Location and qualification\\
\midrule
\endfirsthead
\toprule
Statement & Status & Location and qualification\\
\midrule
\endhead
Inverse coefficient formula at an arbitrary base point & Proved, finite &
Theorem~\ref{thm:inverse-around-a}; exactly the derivative/limit formula after
translating derivatives to coefficients.\\
Analytic inverse has nonzero radius & Proved, named analytic &
Theorem~\ref{thm:quant-inverse}; an explicit rational lower bound is given.\\
Formal-power-series version & Proved, finite &
Theorems~\ref{thm:lb-alt} and \ref{thm:lb-der}.\\
Formula for \(H(G)\) & Proved, finite &
Both displayed Lagrange--B\"urmann forms.\\
Several variables & Proved, finite &
Good's determinant formula, Theorem~\ref{thm:good}.\\
Case \(f'(a)=0\) & Proved, finite ramified form &
Puiseux branches after adjoining the required algebraic roots,
Theorem~\ref{thm:puiseux}.\\
Bell-polynomial coefficient formula & Proved, finite &
Theorem~\ref{thm:bell}.\\
Associahedron face formula & Corrected and proved &
The standard ordinary-coefficient theorem is
Theorem~\ref{thm:assoc}; the page's displayed normalization fails in degree two,
Audit~\ref{audit:assoc}.\\
Solution of \(x^p-x+z=0\) & Proved &
Coefficient formula, boundary convergence, and largest inverse disk in
Section~\ref{sec:algebraic-example}.\\
Lambert \(W\) series at zero & Proved &
Coefficients and exact radius \(e^{-1}\), Section~\ref{sec:lambert}.\\
Shifted Lambert \(W\) series & Proved &
Coefficients, exact radius \(\sqrt{4+\pi^2}\), real interval, and a rational
certificate for \(W(1)\), Section~\ref{sec:shifted-lambert}.\\
Binary trees and Catalan numbers & Proved, finite &
Section~\ref{sec:catalan}; ``unlabelled'' is interpreted as plane rooted full
binary trees, as required by the functional equation.\\
Role in Laplace--Erd\'elyi asymptotics & Made explicit &
The coefficient transformation and finite remainder route are proved in
Section~\ref{sec:laplace}.\\
\bottomrule
\end{longtable}

The page also contains historical attributions and comments about the diversity
of proofs.  Those are source claims rather than mathematical propositions and
are not objects of arithmetic proof here.

\section{Finite series algebra inside \(\ipra\)}\label{sec:finite}

\subsection{Codes and total operations}

Fix a commutative coefficient ring \(R\) whose operations are represented by
primitive-recursive functions.  The principal object-theory instances are
\(\Z\), \(\Q\), Gaussian rationals, and finite algebraic extensions represented
by reduced tuples.  For \(N\in\N\), write
\[
 T_N(R)=R[X]/(X^{N+1}).
\]
We implement an element by its vector \((a_0,\ldots,a_N)\).  Addition and
Cauchy multiplication are
\begin{align*}
 (a+b)_n&=a_n+b_n,\\
 (ab)_n&=\sum_{i=0}^{n}a_i b_{n-i}.
\end{align*}
The derivative and composition, when \(b_0=0\), are
\begin{align*}
 (Da)_n&=(n+1)a_{n+1},\\
 a\circ b&=\sum_{k=0}^{N} a_k b^k \pmod {X^{N+1}}.
\end{align*}
Powers are computed by repeated squaring or primitive recursion.  All displayed
sums are bounded.  Rational division is made total by returning a fixed default
when the denominator is zero; every theorem includes the guard proving that a
denominator actually used is nonzero.

Integers may be coded by sign and magnitude, rationals by a signed numerator
and a positive denominator, and finite vectors by any standard pairing/list
code.  Equality of reduced rational codes and order of rational codes are
decidable.  A finite conjunction \(\forall i<m\,P(i)\) is represented by the
Boolean fold \(\operatorname{all}_{i<m}\chi_P(i)\).  Hence the induction
invariants used below are literally quantifier-free.

\subsection{The constructive fixed point}

Let \(\Phi\in T_N(R)\), with no assumption yet on its constant coefficient.
Define
\[
 \mathcal S_\Phi(A)=X\,\Phi(A),\qquad
 G^{(0)}=0,\qquad G^{(r+1)}=\mathcal S_\Phi(G^{(r)}).
\]
For two truncated series, write \(A\equiv_m B\) when their coefficients below
\(m\) agree.

\begin{lemma}[One-degree improvement]\label{lem:improve}
If \(A\equiv_m B\), then
\[
 \mathcal S_\Phi(A)\equiv_{m+1}\mathcal S_\Phi(B).
\]
\end{lemma}
\begin{proof}
For every \(k\), the coefficient of degree \(j<m\) in \(A^k\) is a finite sum
of products of coefficients of \(A\) whose indices add to \(j\); every index is
therefore below \(m\).  Replacing \(A\) by \(B\) leaves the coefficient
unchanged.  The same holds for the finite sum \(\Phi(A)\).  Multiplication by
\(X\) shifts the agreement bound from \(m\) to \(m+1\).  This is a bounded-sum
argument and uses no case distinction on an undecidable proposition.
\end{proof}

\begin{theorem}[Constructive existence and uniqueness]\label{thm:fixed}
The iterate \(G=G^{(N+1)}\) satisfies
\[
 G=X\Phi(G)\pmod {X^{N+1}}.
\]
It is the unique element of \(T_N(R)\) with zero constant coefficient satisfying
this equation.  The map \((N,\Phi)\mapsto G\) is primitive recursive.
\end{theorem}
\begin{proof}
Lemma~\ref{lem:improve}, applied successively to \(G^{(r)}\) and
\(G^{(r+1)}\), shows
\[
 G^{(r)}\equiv_r G^{(r+1)}.
\]
Thus \(G^{(N+1)}=G^{(N+2)}\) in \(T_N(R)\), which is the fixed-point equation.

If \(A\) and \(B\) are two fixed points with zero constant coefficient, then
\(A\equiv_1B\).  Lemma~\ref{lem:improve} and the two fixed-point equations give
successively \(A\equiv_2B,\ldots,A\equiv_{N+1}B\), hence \(A=B\).  The
construction is an explicitly bounded iteration of explicitly bounded
operations, so it is primitive recursive.
\end{proof}

\begin{remark}
The fixed-point theorem works over an arbitrary commutative ring and does not
require \(\Phi(0)\) to be invertible.  Invertibility enters only when one wants
to rewrite the equation as compositional inversion of \(F(X)=X/\Phi(X)\).
\end{remark}

\section{Finite formal residues}\label{sec:residue}

\subsection{Laurent vectors}

For fixed bounds \(B,N\), a Laurent vector is
\[
 A(X)=\sum_{j=-B}^{N}a_jX^j.
\]
Its residue is \(\Res A=a_{-1}\).  Only a finite interval of exponents is ever
needed in a coefficient proof, so multiplication and substitution are again
bounded algorithms.  We use the elementary identities
\[
 [X^n]A(X)=\Res\bigl(A(X)X^{-n-1}\bigr),
 \qquad \Res(DA)=0.
\]
The second identity is immediate: a derivative can contribute to \(X^{-1}\)
only from differentiating the constant monomial, and its coefficient is zero.

\subsection{Change of variable}

Let
\[
 F(X)=XU(X),\qquad U(0)\in R^\times.
\]
The finite Laurent expression \(Q(F(X))F'(X)\) is well defined to every fixed
requested degree.  The following lemma is the entire residue machinery needed
for Lagrange inversion.

\begin{lemma}[Formal residue substitution]\label{lem:res-subst}
Over a commutative \(\Q\)-algebra,
\[
 \Res\bigl(Q(F(X))F'(X)\bigr)=\Res Q(X)
\]
for every finite Laurent polynomial \(Q\).  The identity is coefficientwise
finite and uniform in all exponent bounds.
\end{lemma}
\begin{proof}
By linearity it suffices to take \(Q(Y)=Y^k\).
If \(k\ne-1\), then
\[
 F(X)^kF'(X)=\frac{1}{k+1}D\bigl(F(X)^{k+1}\bigr),
\]
so its residue is zero, which equals \(\Res Y^k\).
For \(k=-1\),
\[
 \frac{F'}{F}=\frac1X+\frac{U'}{U}.
\]
The inverse of the unit series \(U\) is obtained recursively to the finite
required degree.  The series \(U'/U\) has no negative powers, so the residue is
one.  This is \(\Res Y^{-1}\).

All denominators are fixed nonzero integers.  Equivalently, one may multiply
the identity by their product before arithmetizing it.  Thus the proof is a
finite rational polynomial identity and is intuitionistically valid.
\end{proof}

\begin{remark}[Integral version]
The alternative Lagrange identity below has integral coefficients after
clearing only unit inverses.  Therefore it can first be proved as a universal
integer polynomial identity in finitely many coefficient variables and then
interpreted in any desired commutative ring.  The \(\Q\)-algebra assumption in
Lemma~\ref{lem:res-subst} is a convenient proof environment, not an essential
source of nonconstructivity.
\end{remark}

\section{The constructive Lagrange--B\"urmann theorem}\label{sec:main}

Let \(\Phi,H\in\fps{R}\), where \(R\) is a commutative \(\Q\)-algebra, and
let \(G\) be the unique zero-constant solution
\begin{equation}\label{eq:fixed}
 G(Z)=Z\Phi(G(Z)).
\end{equation}
For each requested coefficient, this means the finite solution of
Theorem~\ref{thm:fixed}.

Assume first that \(\Phi(0)\) is a unit and put
\[
 F(W)=\frac{W}{\Phi(W)}.
\]
Then \(F(G(Z))=Z\) and \(G(F(W))=W\), coefficientwise.  The first identity is
just \eqref{eq:fixed}; the second follows either from finite uniqueness or by
applying the same recursive inverse construction.

\begin{theorem}[Alternative Lagrange--B\"urmann formula]\label{thm:lb-alt}
For \(n\ge1\),
\begin{equation}\label{eq:lb-alt}
 [Z^n]H(G(Z))
 = [W^n]H(W)\Phi(W)^{n-1}
       \bigl(\Phi(W)-W\Phi'(W)\bigr).
\end{equation}
The assertion for coefficient \(n\) depends only on coefficients through
finite degree \(n\), and is formalizable uniformly in \(n\) in \(\ipra\).
\end{theorem}
\begin{proof}
Use coefficient-as-residue and substitute \(Z=F(W)\) by
Lemma~\ref{lem:res-subst}:
\begin{align*}
 [Z^n]H(G(Z))
 &=\Res_Z H(G(Z))Z^{-n-1}\\
 &=\Res_W H(G(F(W)))F(W)^{-n-1}F'(W)\\
 &=\Res_W H(W)
   \left(\frac{\Phi(W)}{W}\right)^{n+1}
   \frac{\Phi(W)-W\Phi'(W)}{\Phi(W)^2}\\
 &=\Res_W H(W)W^{-n-1}\Phi(W)^{n-1}
       \bigl(\Phi(W)-W\Phi'(W)\bigr).
\end{align*}
The last residue is exactly the coefficient on the right of
\eqref{eq:lb-alt}.  Although unit notation made the calculation compact, after
multiplying out, the right side contains no negative power of \(\Phi\).  Hence
both sides are universal polynomials in the finitely many coefficients involved.
The identity therefore extends from the unit case to arbitrary \(\Phi\) by the
standard polynomial-identity argument, or directly by adjoining an
indeterminate unit and comparing coefficients.  This extension is again
finite and constructive.
\end{proof}

\begin{theorem}[Derivative Lagrange--B\"urmann formula]\label{thm:lb-der}
For \(n\ge1\),
\begin{equation}\label{eq:lb-der-cleared}
 n[Z^n]H(G(Z))=[W^{n-1}]H'(W)\Phi(W)^n.
\end{equation}
Over \(\Q\), equivalently,
\begin{equation}\label{eq:lb-der}
 [Z^n]H(G(Z))=\frac1n[W^{n-1}]H'(W)\Phi(W)^n.
\end{equation}
\end{theorem}
\begin{proof}
Differentiate the finite product \(H\Phi^n\).  Since
\([W^{n-1}](A')=n[W^n]A\),
\begin{align*}
 [W^{n-1}]H'\Phi^n
 &=n[W^n]H\Phi^n
   -n[W^{n-1}]H\Phi^{n-1}\Phi'\\
 &=n[W^n]H\Phi^{n-1}(\Phi-W\Phi').
\end{align*}
Theorem~\ref{thm:lb-alt} identifies the last coefficient with
\(n[Z^n]H(G)\).  This proof uses only the Leibniz rule and finite coefficient
arithmetic.  Formula \eqref{eq:lb-der-cleared}, rather than division by \(n\),
is the preferred object-theory statement.
\end{proof}

\begin{corollary}[Basic Lagrange inversion]\label{cor:basic}
For \(n\ge1\),
\begin{equation}\label{eq:basic}
 [Z^n]G(Z)=\frac1n[W^{n-1}]\Phi(W)^n.
\end{equation}
More generally, for \(k\ge1\),
\begin{equation}\label{eq:power}
 [Z^n]G(Z)^k=\frac{k}{n}[W^{n-k}]\Phi(W)^n,
\end{equation}
where the right side is zero for \(n<k\).
\end{corollary}
\begin{proof}
Take \(H(W)=W\) and \(H(W)=W^k\) in
Theorem~\ref{thm:lb-der}.  In the second case
\(H'=kW^{k-1}\), and shifting the extracted exponent gives
\eqref{eq:power}.
\end{proof}

\subsection{The induction proof encoded by the finite construction}

The residue proof is conceptually short, but the fixed-point construction also
gives a direct quantifier-free induction route.  Define the candidate
coefficient
\[
 L_n(\Phi)=\frac1n[W^{n-1}]\Phi(W)^n.
\]
The recurrence obtained by comparing coefficients in \(G=Z\Phi(G)\) expresses
\(g_n\) as a bounded polynomial in \(\phi_0,\ldots,\phi_{n-1}\) and
\(g_1,\ldots,g_{n-1}\).  Substitution of \(L_1,\ldots,L_{n-1}\) reduces the
step to the finite derivative identity used above.  The invariant
``all coefficients below \(r\) agree with \(L_i\)'' is a Boolean fold, so
ProveIt's \code{PRA.qfInd} applies without adding bounded quantifiers to the
induction formula.  This gives an alternative object-level proof if one prefers
to avoid formal Laurent syntax.

\section{Inverse series around an arbitrary point}\label{sec:inverse}

Let \(f\) be a formal or analytic germ at \(a\), and put
\[
 F(X)=f(a+X)-f(a).
\]
Suppose \(F'(0)=f'(a)\) is a unit.  Define
\[
 \Phi(X)=\frac{X}{F(X)}.
\]
Let \(G\) solve \(G=Z\Phi(G)\).  Then
\[
 F(G(Z))=Z,
\]
so the inverse germ is
\[
 g(f(a)+Z)=a+G(Z).
\]

\begin{theorem}[Inverse coefficient formula at \(a\)]\label{thm:inverse-around-a}
Write the inverse in exponential-coefficient convention as
\[
 g(z)=a+\sum_{n\ge1}g_n\frac{(z-f(a))^n}{n!}.
\]
Then
\begin{equation}\label{eq:inverse-a-coeff}
 g_n=(n-1)![X^{n-1}]
 \left(\frac{X}{f(a+X)-f(a)}\right)^n.
\end{equation}
If derivatives are available, this is the page's formula
\begin{equation}\label{eq:inverse-a-deriv}
 g_n=\left.
 \frac{d^{n-1}}{dw^{n-1}}
 \left(\frac{w-a}{f(w)-f(a)}\right)^n
 \right|_{w=a}.
\end{equation}
\end{theorem}
\begin{proof}
Corollary~\ref{cor:basic} gives
\([Z^n]G=\frac1n[X^{n-1}]\Phi^n\).  Since
\([Z^n]G=g_n/n!\), multiplication by \(n!\) gives
\eqref{eq:inverse-a-coeff}.  For a formal power series, the coefficient--
derivative relation
\[
 A^{(m)}(0)=m![X^m]A
\]
is proved by \(m\) finite differentiations.  Substituting \(X=w-a\) gives
\eqref{eq:inverse-a-deriv}.  The ``limit'' notation on the page merely denotes
evaluation at a removable singularity: the quotient has unit constant term
\(1/f'(a)\).
\end{proof}

\section{Explicit coefficient formulae}\label{sec:explicit}

\subsection{Bell polynomials}

For \(m,k\in\N\), define the exponential partial Bell polynomial by the finite
sum
\begin{equation}\label{eq:bell-def}
 B_{m,k}(x_1,x_2,\ldots)
 =\sum_{\substack{r_1+r_2+\cdots=k\\
                   r_1+2r_2+\cdots=m}}
 \frac{m!}{r_1!r_2!\cdots}
 \prod_{j\ge1}\left(\frac{x_j}{j!}\right)^{r_j}.
\end{equation}
Only \(j\le m-k+1\) occurs.  The sum is a bounded search over multiplicity
vectors and is primitive recursive.  Its defining generating identity is
\begin{equation}\label{eq:bell-power}
 \left(\sum_{j\ge1}x_j\frac{W^j}{j!}\right)^k
 =k!\sum_{m\ge k}B_{m,k}(x_1,x_2,\ldots)\frac{W^m}{m!}.
\end{equation}
It follows directly by the finite multinomial theorem.

\begin{theorem}[Bell-polynomial inverse coefficients]\label{thm:bell}
Suppose
\[
 f(W)=\sum_{j\ge0}f_j\frac{W^j}{j!},\qquad
 g(Z)=\sum_{j\ge0}g_j\frac{Z^j}{j!},
\]
with \(f_0=0\), \(f_1\ne0\), and \(f(g(Z))=Z\).  Put
\[
 \widehat f_j=\frac{f_{j+1}}{(j+1)f_1},
 \qquad n^{\overline{k}}=n(n+1)\cdots(n+k-1).
\]
Then \(g_1=f_1^{-1}\), and for \(n\ge2\),
\begin{equation}\label{eq:bell-result}
 g_n=f_1^{-n}\sum_{k=1}^{n-1}(-1)^k n^{\overline{k}}
 B_{n-1,k}(\widehat f_1,\ldots,\widehat f_{n-k}).
\end{equation}
\end{theorem}
\begin{proof}
Factor
\[
 f(W)=f_1W(1+A(W)),\qquad
 A(W)=\sum_{j\ge1}\widehat f_j\frac{W^j}{j!}.
\]
Thus \(\Phi(W)=W/f(W)=f_1^{-1}(1+A)^{-1}\).  By
Theorem~\ref{thm:inverse-around-a},
\[
 g_n=(n-1)!f_1^{-n}[W^{n-1}](1+A)^{-n}.
\]
To degree \(n-1\), the negative binomial expansion is the finite identity
\[
 (1+A)^{-n}
 =\sum_{k=0}^{n-1}(-1)^k\frac{n^{\overline{k}}}{k!}A^k.
\]
For \(n\ge2\), the \(k=0\) term has no coefficient of degree \(n-1\).
Equation~\eqref{eq:bell-power} gives
\[
 [W^{n-1}]A^k=\frac{k!}{(n-1)!}B_{n-1,k}.
\]
Substitution and cancellation yield \eqref{eq:bell-result}.  Every expansion is
truncated before use, so no convergence or infinitary binomial theorem is
invoked.
\end{proof}

\subsection{Associahedra and the normalization correction}\label{sec:assoc}

Use the convention of Aguiar--Ardila: the associahedron \(\mathfrak a_n\) has
dimension \(n-1\), and each face is combinatorially a product
\[
 F\cong \mathfrak a_{i_1}\times\cdots\times\mathfrak a_{i_r}.
\]
For
\[
 C(X)=X+\sum_{n\ge1}c_nX^{n+1},\qquad
 D(X)=X+\sum_{n\ge1}d_nX^{n+1},
 \qquad C(D(X))=X,
\]
write \(c_F=c_{i_1}\cdots c_{i_r}\).

\begin{theorem}[Polytopal compositional inversion]\label{thm:assoc}
For \(n\ge1\),
\begin{equation}\label{eq:assoc-standard}
 d_n=\sum_{F\text{ face of }\mathfrak a_n}
       (-1)^{n-\dim F}c_F.
\end{equation}
\end{theorem}
\begin{proof}
Faces of \(\mathfrak a_n\) are in bijection with plane rooted trees having
\(n+1\) leaves and internal vertices of arity at least two.  A vertex of arity
\(j+1\) contributes the factor \(c_j\).  If the tree has \(r\) internal
vertices, the corresponding face has dimension \(n-r\), hence sign
\((-1)^r=(-1)^{n-\dim F}\).

Now expand the identity \(C(D)=X\).  For positive degree it is
\[
 D-X=-\sum_{j\ge1}c_jD^{j+1}.
\]
Repeatedly replacing every occurrence of \(D\) on the right by the same
expression terminates at each fixed degree.  A term is exactly a plane rooted
tree; every internal vertex contributes one minus sign and the coefficient
\(c_j\) determined by its arity.  Therefore the coefficient of \(X^{n+1}\) is
the signed tree sum, and the face--tree bijection turns it into
\eqref{eq:assoc-standard}.  This is a finite expansion at each degree and is
formalizable by induction on the number of leaves.
\end{proof}

\begin{audit}[The page's displayed associahedron formula is not normalized
consistently]\label{audit:assoc}
Immediately before its face formula, the page uses exponential coefficients
\(f(W)=\sum f_nW^n/n!\) and \(g(Z)=\sum g_nZ^n/n!\), with \(f_1=1\).
The standard theorem applies instead to ordinary coefficients
\[
 c_m=\frac{f_{m+1}}{(m+1)!},\qquad
 d_m=\frac{g_{m+1}}{(m+1)!}.
\]
Consequently the correct translation is
\begin{equation}\label{eq:assoc-exp-correct}
 g_n=n!\sum_{F\text{ face of }\mathfrak a_{n-1}}
 (-1)^{n-1-\dim F}
 \prod_{\mathfrak a_j\text{ factor of }F}
       \frac{f_{j+1}}{(j+1)!},\qquad n\ge2.
\end{equation}
For the degree-two test
\(f(W)=W+cW^2/2\), direct inversion gives
\(g(Z)=Z-cZ^2/2+\cdots\), hence \(g_2=-c\).  Formula
\eqref{eq:assoc-exp-correct} gives exactly this.  The page's displayed product
of unshifted, unscaled \(f_i\)'s does not.  Thus one cannot prove that display
as written; the correct ordinary formula and its exponential translation are
what is proved here.
\end{audit}

\section{Multivariate and ramified extensions}\label{sec:extensions}

\subsection{Good's multivariate formula}

Let \(d\ge1\), let \(\mathbf Z=(Z_1,\ldots,Z_d)\), and suppose the
multivariate series \(G_i(\mathbf Z)\) are defined by
\begin{equation}\label{eq:multi-fixed}
 G_i=Z_i\Phi_i(G_1,\ldots,G_d),\qquad 1\le i\le d,
\end{equation}
where each \(\Phi_i(0)\) is a unit.  Iteration from the zero tuple improves
total-degree agreement by one, exactly as in Lemma~\ref{lem:improve}; hence the
solution exists uniquely at every finite truncation order.

For a multi-index \(\mathbf n=(n_1,\ldots,n_d)\), use
\(\mathbf Z^{\mathbf n}=\prod_iZ_i^{n_i}\).  Define the logarithmic Jacobian
matrix
\begin{equation}\label{eq:good-matrix}
 L_{ij}(\mathbf W)=
 \delta_{ij}-\frac{W_j}{\Phi_i(\mathbf W)}
                    \frac{\partial\Phi_i}{\partial W_j}(\mathbf W).
\end{equation}

\begin{theorem}[Good's multivariate Lagrange formula]\label{thm:good}
For every multivariate series \(H\),
\begin{equation}\label{eq:good}
 [\mathbf Z^{\mathbf n}]H(\mathbf G(\mathbf Z))
 = [\mathbf W^{\mathbf n}]
 H(\mathbf W)\prod_{i=1}^d\Phi_i(\mathbf W)^{n_i}
 \det L(\mathbf W).
\end{equation}
For fixed requested degree and dimension, the formula is a finite identity.  A
uniform-in-\(d\) version is also primitive recursive when vectors and matrices
are coded as finite lists and the determinant is the bounded Leibniz sum.
\end{theorem}
\begin{proof}
Put \(F_i(\mathbf W)=W_i/\Phi_i(\mathbf W)\).  Multivariate coefficient
extraction is an iterated finite residue:
\[
 [\mathbf Z^{\mathbf n}]A
 =\Res_{Z_1}\cdots\Res_{Z_d}
 A(\mathbf Z)\prod_iZ_i^{-n_i-1}.
\]
The multivariate finite change-of-variables lemma is proved on Laurent
monomials, one variable at a time, or by expanding the determinant and applying
Lemma~\ref{lem:res-subst}.  Substituting \(\mathbf Z=\mathbf F(\mathbf W)\)
gives \(\mathbf G(\mathbf F(\mathbf W))=\mathbf W\) and contributes the
ordinary Jacobian determinant.

A direct derivative calculation yields
\[
 \frac{\partial F_i}{\partial W_j}
 =\frac1{\Phi_i}
   \left(\delta_{ij}-\frac{W_i}{\Phi_i}
                     \frac{\partial\Phi_i}{\partial W_j}\right).
\]
After multiplying by \(\prod_iF_i^{-n_i-1}\), the row factors
\(\Phi_i^{-1}\) cancel one power from each \(\Phi_i^{n_i+1}\).  The remaining
determinant may be written with row factors \(W_i/\Phi_i\), or in the
logarithmic form \eqref{eq:good-matrix}.  The equality of the two determinants
is \(\det(I-AB)=\det(I-BA)\), itself a finite polynomial identity.  The residue
left after cancellation is exactly the coefficient on the right of
\eqref{eq:good}.
\end{proof}

\subsection{Vanishing first derivative and Puiseux branches}

Suppose, after translating the base point and value to zero, that
\begin{equation}\label{eq:ramified-F}
 F(X)=cX^mV(X),\qquad m\ge2,
 \qquad c\in R^\times,\quad V(0)=1.
\end{equation}
Work in a finite algebraic extension containing a chosen \(m\)-th root
\(\beta\) of \(c^{-1}\), all needed \(m\)-th roots of unity \(\zeta_j\), and
\(1/m\).  The binomial series
\[
 V(X)^{-1/m}=\sum_{r\ge0}\binom{-1/m}{r}(V(X)-1)^r
\]
is finite to every requested degree because \(V-1\) has zero constant term.
Set
\[
 \Phi_j(X)=\zeta_j\beta V(X)^{-1/m}.
\]

\begin{theorem}[Constructive ramified inversion]\label{thm:puiseux}
For each \(m\)-th root of unity \(\zeta_j\), there is a unique formal series
\(X_j(T)\) with zero constant term satisfying
\[
 X_j(T)=T\Phi_j(X_j(T)).
\]
It obeys
\[
 F(X_j(T))=T^m,
\]
and its coefficients are
\begin{equation}\label{eq:puiseux-coeff}
 [T^n]X_j(T)=\frac1n[X^{n-1}]\Phi_j(X)^n.
\end{equation}
Consequently
\[
 a+X_j\bigl((z-f(a))^{1/m}\bigr)
\]
are the formal Puiseux branches of the multivalued inverse near \(a\).
\end{theorem}
\begin{proof}
Existence, uniqueness, and \eqref{eq:puiseux-coeff} are
Theorem~\ref{thm:fixed} and Corollary~\ref{cor:basic}.  Raising the fixed-point
equation to the \(m\)-th power gives
\[
 X_j(T)^m=T^m\zeta_j^m\beta^mV(X_j(T))^{-1}
 =T^m c^{-1}V(X_j(T))^{-1}.
\]
Multiplication by \(cV(X_j(T))\) proves \(F(X_j(T))=T^m\).  Every operation in
this argument is finite in a truncation, including construction of the
binomial root.  Analytic branches require the corresponding effective choices
of algebraic roots and a slit neighbourhood, but no new coefficient identity.
\end{proof}

\section{Quantitative analytic lifting}\label{sec:analytic}

The formal theorem proves the Taylor coefficients of an inverse if an analytic
inverse exists.  We now construct the inverse with an explicit positive radius.
The constants are intentionally simple rather than optimal.

Translate so that the base point and value are zero.  Suppose
\begin{equation}\label{eq:analytic-F}
 F(w)=a_1w+\sum_{n\ge2}a_nw^n
\end{equation}
comes with positive rationals \(R,M,K\) satisfying
\begin{equation}\label{eq:analytic-bounds}
 |a_n|\le MR^{-n}\quad(n\ge2),
 \qquad |a_1^{-1}|\le K.
\end{equation}
The coefficients, the inequalities, and the tail modulus are assumed to be
primitive-recursively named.  When \(M>0\), set
\begin{equation}\label{eq:radii}
 t=\min\left(\frac12,\frac{R}{16KM}\right),
 \qquad \rho=Rt,
 \qquad r=\frac{\rho}{2K}.
\end{equation}
When \(M=0\), take \(t=1/2\).  All numbers are positive rationals.

\begin{lemma}[Derivative control]\label{lem:derivative-control}
Writing \(h(w)=F(w)-a_1w\), one has on \(|w|\le\rho\)
\[
 |h'(w)|\le\frac1{2K}.
\]
\end{lemma}
\begin{proof}
For \(0\le t\le1/2\), finite geometric differentiation followed by the tail
bound gives
\[
 \sum_{n\ge2}nt^{n-1}=\frac1{(1-t)^2}-1\le8t.
\]
Therefore
\[
 |h'(w)|
 \le\frac{M}{R}\sum_{n\ge2}nt^{n-1}
 \le\frac{8Mt}{R}\le\frac1{2K}.
\]
Each infinite-looking inequality means a rational partial-sum inequality plus
a geometric remainder estimate with an explicit cutoff.  It is therefore an
\(\ipra\) theorem for the supplied coefficient name.
\end{proof}

\begin{theorem}[Explicit constructive inverse radius]\label{thm:quant-inverse}
For every \(|z|\le r\), the iteration
\begin{equation}\label{eq:picard}
 w_0(z)=0,
 \qquad w_{q+1}(z)=a_1^{-1}\bigl(z-h(w_q(z))\bigr)
\end{equation}
converges uniformly to a unique \(g(z)\) with \(|g(z)|\le\rho\) and
\(F(g(z))=z\).  Moreover
\begin{equation}\label{eq:picard-error}
 |w_q(z)-g(z)|\le\rho\,2^{-q}.
\end{equation}
The function \(g\) is analytic on \(|z|<r\), and its Taylor coefficients are
those in Theorem~\ref{thm:inverse-around-a}.  Thus the inverse series has a
nonzero, explicitly bounded radius of convergence.
\end{theorem}
\begin{proof}
For \(|w|\le\rho\), Lemma~\ref{lem:derivative-control} and the segment integral
estimate give \(|h(w)|\le\rho/(2K)\).  Hence, if \(|z|\le r\),
\[
 \left|a_1^{-1}(z-h(w))\right|
 \le K r+K\frac{\rho}{2K}=\rho.
\]
Thus the closed radius-\(\rho\) ball maps into itself.  For two points in the
ball,
\[
 \left|a_1^{-1}(h(u)-h(v))\right|\le\frac12|u-v|,
\]
so the map is a contraction with the displayed rational constant.  Since
\(|w_1-w_0|\le Kr=\rho/2\), the geometric tail of successive differences gives
\eqref{eq:picard-error}.  This directly constructs a regular name for \(g(z)\)
and proves uniqueness.

Every iterate \(w_q(z)\) is analytic in \(z\).  On any rational closed subdisk
of \(|z|<r\), the error bound is uniform and primitive recursive, so termwise
Cauchy estimates, or the existing constructive Cauchy package, show that the
limit is analytic with a primitive-recursive modulus.  Passing to the limit in
\eqref{eq:picard} gives \(F(g)=z\).

Finally, the coefficient vector of \(w_q\) through degree \(N\) stabilizes once
\(q>N\), by Lemma~\ref{lem:improve}.  It is therefore the finite formal fixed
point, so Corollary~\ref{cor:basic} and
Theorem~\ref{thm:inverse-around-a} give the Taylor coefficients of the analytic
limit.  No appeal to excluded middle, compactness, or an unspecified inverse
function theorem is made.
\end{proof}

\begin{remark}[Injectivity]
The same estimate proves directly that, for \(u,v\) in the radius-\(\rho\)
ball,
\[
 |F(u)-F(v)|\ge |a_1||u-v|-|h(u)-h(v)|
 \ge\frac1{2K}|u-v|.
\]
Thus the local inverse is genuinely single-valued on the constructed disk.
\end{remark}

\section{The algebraic equation \(x^p-x+z=0\)}\label{sec:algebraic-example}

Fix an integer \(p\ge2\).  The branch with \(x(0)=0\) satisfies
\[
 x=z(1-x^{p-1})^{-1}.
\]
Thus \(\Phi(w)=(1-w^{p-1})^{-1}\).  By
Corollary~\ref{cor:basic},
\[
 [z^n]x(z)=\frac1n[w^{n-1}](1-w^{p-1})^{-n}.
\]
The finite negative-binomial identity shows that this coefficient vanishes
unless \(n=(p-1)k+1\).  In that case
\begin{align*}
 [z^{(p-1)k+1}]x(z)
 &=\frac1{(p-1)k+1}
   \binom{(p-1)k+1+k-1}{k}\\
 &=\frac1{(p-1)k+1}\binom{pk}{k}.
\end{align*}
Hence
\begin{equation}\label{eq:algebraic-series}
 x(z)=\sum_{k\ge0}\binom{pk}{k}
       \frac{z^{(p-1)k+1}}{(p-1)k+1},
\end{equation}
which is the displayed series on the page.

\subsection{Exact convergence disk}

Put
\[
 q_p=\frac{(p-1)^{p-1}}{p^p},
 \qquad \rho_p=q_p^{1/(p-1)}=(p-1)p^{-p/(p-1)}.
\]
For
\(a_k=\binom{pk}{k}/((p-1)k+1)\), direct cancellation of factorials gives
\[
 \lim_{k\to\infty}\frac{a_{k+1}}{a_k}
 =\frac{p^p}{(p-1)^{p-1}}=q_p^{-1}.
\]
The limit is constructive: after writing the ratio as a product of the fixed
number \(p\) of rational factors, each factor differs from its limit by at most
\(C_p/(k+1)\).  Therefore the series in the variable
\(t=z^{p-1}\) has radius \(q_p\), and \eqref{eq:algebraic-series} has radius
\(\rho_p\) in \(z\).

The page includes the boundary.  A finite Wallis/Stirling estimate gives, for
a rational constant \(C_p\) depending only on the fixed integer \(p\),
\begin{equation}\label{eq:alg-boundary}
 \binom{pk}{k}q_p^k\le\frac{C_p}{\sqrt{k+1}}.
\end{equation}
One constructive derivation expands the factorial quotient into finite products,
applies the rational inequalities
\(u-u^2/2\le\log(1+u)\le u\) for \(0\le u\le1\), and bounds the resulting
harmonic sums by dyadic blocks.  Squaring clears the square root, and all log
bounds are finite Taylor estimates with rational remainders.  Thus
\[
 a_kq_p^k\le\frac{C'_p}{(k+1)^{3/2}},
\]
and comparison with the telescoping dyadic grouping of
\(\sum(k+1)^{-3/2}\) proves absolute convergence for \(|z|=\rho_p\).  This
proof has a primitive-recursive cutoff for every requested error.

Finally, for \(f(x)=x-x^p\),
\[
 f'(x)=1-px^{p-1}.
\]
Its critical points satisfy \(x^{p-1}=1/p\), and their critical values are
\[
 f(x)=\frac{p-1}{p}x,
 \qquad |f(x)|=(p-1)p^{-p/(p-1)}=\rho_p.
\]
The branch point at each such value prevents analytic continuation of the local
inverse through the corresponding point; conversely the convergent series
constructs the inverse throughout the open disk.  Therefore this is the largest
centered disk on which the local branch at zero is analytic.

\section{Lambert \(W\)}\label{sec:lambert}

The principal germ is defined by
\[
 W(z)e^{W(z)}=z,
\]
or equivalently
\[
 W(z)=ze^{-W(z)}.
\]
Thus \(\Phi(w)=e^{-w}\).  Corollary~\ref{cor:basic} gives
\begin{align*}
 [z^n]W(z)
 &=\frac1n[w^{n-1}]e^{-nw}
 =\frac1n\frac{(-n)^{n-1}}{(n-1)!}\\
 &=\frac{(-n)^{n-1}}{n!}.
\end{align*}
Consequently
\begin{equation}\label{eq:W-series}
 W(z)=\sum_{n\ge1}\frac{(-n)^{n-1}}{n!}z^n
 =z-z^2+\frac32z^3-\frac83z^4+\frac{125}{24}z^5-\cdots.
\end{equation}
Formal exponentiation to degree \(N\) uses only the rational coefficients
\(1/j!\) for \(j\le N\), so the calculation is finite in \(\ipra\).

For absolute coefficients \(c_n=n^{n-1}/n!\),
\begin{equation}\label{eq:W-ratio}
 \frac{c_{n+1}}{c_n}=\left(1+\frac1n\right)^{n-1}.
\end{equation}
The elementary binomial bounds defining the primitive-recursive name of \(e\)
show that the right side tends to \(e\).  The radius of
\eqref{eq:W-series} is therefore \(e^{-1}\).  At \(w=-1\), the map
\(w\mapsto we^w\) has derivative zero and critical value \(-e^{-1}\), so the
radius cannot be enlarged.  The branch obtained from the germ at zero is, by
definition and by real monotonicity on \((-1,\infty)\), the principal branch.

The same finite factorial estimate used in
\eqref{eq:alg-boundary} gives
\(c_ne^{-n}=O(n^{-3/2})\), so the series also converges absolutely at every
point of the boundary circle \(\lvert z\rvert=e^{-1}\), since taking absolute
values removes the phase.  This strengthens, but is not needed for, the page's radius claim.

\section{The shifted Lambert series}\label{sec:shifted-lambert}

Put
\[
 U(z)=W(e^{1+z})-1.
\]
Since \((1+U)e^{1+U}=e^{1+z}\), taking the local logarithm at \(U=0\) gives
\begin{equation}\label{eq:shift-equation}
 z=U+\log(1+U)=:A(U).
\end{equation}
Here
\[
 A(U)=2U-\frac{U^2}{2}+\frac{U^3}{3}-\frac{U^4}{4}+\cdots.
\]
Applying Theorem~\ref{thm:inverse-around-a}, or simply the finite inverse
recursion, yields
\begin{equation}\label{eq:shift-series}
 \begin{split}
 W(e^{1+z})={}&1+\frac z2+\frac{z^2}{16}-\frac{z^3}{192}
 -\frac{z^4}{3072}+\frac{13z^5}{61440}\\
 &-\frac{47z^6}{1474560}+O(z^7).
 \end{split}
\end{equation}
The archive's exact verifier reproduces these rational coefficients.

\subsection{Exact radius \(\sqrt{4+\pi^2}\)}

A convenient global parameter is
\[
 s=\log(1+U),\qquad U=e^s-1.
\]
Then \eqref{eq:shift-equation} becomes inversion of the entire function
\begin{equation}\label{eq:E}
 E(s)=e^s+s-1.
\end{equation}
Its critical points and values are
\begin{align*}
 E'(s)=e^s+1=0
 &\iff s=(2k+1)\pi i,\\
 E((2k+1)\pi i)
 &=-2+(2k+1)\pi i.
\end{align*}
The nearest critical values to zero therefore have modulus
\begin{equation}\label{eq:rshift}
 R_*=\sqrt{4+\pi^2}.
\end{equation}

For completeness, there is no missing finite asymptotic obstruction.  Suppose
a continuous curve \(s(t)=x(t)+iy(t)\) tends to infinity while \(E(s(t))\)
remains bounded.  From
\[
 e^{s(t)}=1+E(s(t))-s(t)
\]
one obtains \(e^{x(t)}\asymp |s(t)|\), hence
\(x(t)=\log|s(t)|+O(1)\), \(|y(t)|\to\infty\), and
\(x(t)/|y(t)|\to0\).  Dividing the displayed equation by its magnitude gives
\[
 e^{iy(t)}\longrightarrow -i\,\operatorname{sgn}(y(t)).
\]
For sufficiently large \(t\), the continuous value \(y(t)\) must therefore lie
in a union of disjoint short intervals around
\(3\pi/2+2\pi k\) when \(y>0\), or around \(\pi/2+2\pi k\) when \(y<0\).
A connected tail of a continuous curve cannot move between these separated
intervals, so \(y(t)\) would be bounded, a contradiction.  Thus \(E\) has no
finite asymptotic value.

It follows that over a disk avoiding the critical values, each component of the
inverse surface is an unramified covering.  The disk \(|z|<R_*\) is simply
connected, so the component containing \(s=0\) gives a single analytic inverse.
At \(s=\pm\pi i\),
\[
 E''(s)=e^s=-1\ne0,
\]
so the inverse has a square-root branch point at the corresponding nearest
critical value.  Hence the Taylor radius is exactly \(R_*\).

This argument has a constructive arithmetic reading.  For each rational
\(r<R_*\), cover the closed disk \(|z|\le r\) by finitely many rational boxes,
use the explicit lower bound on \(|E'|\) over the finitely traced inverse branch,
and continue by contraction from box to box.  The no-asymptotic estimate gives
a computable bound on \(|s|\).  At the critical value, the nonzero quadratic
term gives a finite certificate that no analytic inverse can cross.  Thus the
claims ``converges below \(R_*\)'' and ``cannot have radius above \(R_*\)'' are
families of first-order \(\ipra\) statements with primitive-recursive moduli.

For positive real \(x\), substitute \(z=\log x-1\).  The series converges when
\begin{equation}\label{eq:W-real-interval}
 |\log x-1|<R_*,
 \qquad	ext{equivalently}
 \qquad e^{1-R_*}<x<e^{1+R_*}.
\end{equation}
Numerically this is
\[
 0.0655991999994\ldots<x<112.639423941\ldots,
\]
consistent with the interval quoted on the page.

\subsection{A fully rational certificate for \(W(1)\)}

The function \(u\mapsto ue^u\) is strictly increasing for \(u\ge0\).  Let
\[
 L=\frac{5671432904097838}{10^{16}},\qquad
 U=\frac{5671432904097840}{10^{16}}.
\]
For \(0\le x<1\), put
\[
 S_{30}(x)=\sum_{j=0}^{30}\frac{x^j}{j!}.
\]
The positive tail satisfies the rational bound
\[
 S_{30}(x)\le e^x\le
 S_{30}(x)+\frac{x^{31}/31!}{1-x/32}.
\]
Exact integer arithmetic verifies
\[
 L\left(S_{30}(L)+\frac{L^{31}/31!}{1-L/32}\right)<1,
 \qquad U S_{30}(U)>1.
\]
Therefore
\begin{equation}\label{eq:omega-cert}
 0.5671432904097838<W(1)<0.5671432904097840.
\end{equation}
This is stronger than the six-decimal approximation on the page and is checked
by \code{code/verify.py} without floating-point arithmetic.

\section{Binary trees and Catalan numbers}\label{sec:catalan}

The functional equation on the page counts plane rooted full binary trees with
unlabelled nodes: a non-leaf consists of an ordered left subtree and an ordered
right subtree.  Without the plane/order structure, the count would instead be
the Wedderburn--Etherington sequence and the displayed product equation would
not apply.

Let \(B_n\) count these trees by internal nodes and
\(B(z)=\sum_{n\ge0}B_nz^n\).  The root decomposition is a constructive
bijection
\[
 \B\cong\{\text{leaf}\}\;\sqcup\;
 \{\text{root}\}\times\B\times\B,
\]
which gives
\[
 B(z)=1+zB(z)^2.
\]
Set \(C=B-1\).  Then
\[
 C=z(1+C)^2.
\]
Corollary~\ref{cor:basic}, with \(\Phi(w)=(1+w)^2\), yields for \(n\ge1\)
\begin{align*}
 B_n=[z^n]C(z)
 &=\frac1n[w^{n-1}](1+w)^{2n}\\
 &=\frac1n\binom{2n}{n-1}
 =\frac1{n+1}\binom{2n}{n}.
\end{align*}
This is the \(n\)-th Catalan number.  The divisibility implicit in the last
fraction follows already from the tree bijection, or algebraically from
\(\binom{2n}{n-1}=n\binom{2n}{n}/(n+1)\).  Finite trees can be coded by balanced
parenthesis words, and the root decomposition and size calculation are
primitive recursive, so both the generating-function and bijective proofs fit
\(\ipra\).

\section{The inversion step in Laplace--Erd\'elyi asymptotics}\label{sec:laplace}

Consider, first in the regular case,
\[
 I(\lambda)=\int_0^b e^{-\lambda f(t)}h(t)\,dt,
 \qquad f(0)=0,\quad f'(0)\ne0.
\]
Let \(u=f(t)\), \(t=g(u)\), and write \(f(t)=t/\Phi(t)\).  If \(K'(t)=h(t)\)
and \(K(0)=0\), then the transformed amplitude is
\[
 A(u)=h(g(u))g'(u)=\frac{d}{du}K(g(u)).
\]
Theorem~\ref{thm:lb-der} gives, for \(m\ge0\),
\begin{equation}\label{eq:laplace-coeff}
 [u^m]A(u)
 =(m+1)[u^{m+1}]K(g(u))
 =[t^m]h(t)\Phi(t)^{m+1}.
\end{equation}
Thus the crucial inversion step produces the amplitude coefficients directly,
without first computing all coefficients of \(g\).

If
\[
 A(u)=\sum_{m=0}^{M-1}a_mu^m+u^MR_M(u)
\]
with a supplied bound \(|R_M(u)|\le C\) on \([0,\delta]\), then splitting the
integral at \(\delta\) and integrating the finite polynomial gives
\begin{equation}\label{eq:laplace-asymp}
 I(\lambda)
 =\sum_{m=0}^{M-1}\frac{m!a_m}{\lambda^{m+1}}
 +\mathcal R_M(\lambda),
\end{equation}
where
\[
 |\mathcal R_M(\lambda)|
 \le \frac{C M!}{\lambda^{M+1}}
 +e^{-\lambda\delta}\int_\delta^{f(b)}|A(u)|\,du.
\]
With rational polygonal integration and the moduli used in the ProveIt
Taylor/Cauchy development, this is a primitive-recursive error bound.  If the
leading behavior is \(f(t)=ct^mV(t)\), use the ramified inverse from
Theorem~\ref{thm:puiseux}; the resulting fractional powers are precisely the
Erd\'elyi generalization.  This makes mathematically explicit the page's
statement that function inversion is the crucial step.

\section{Arithmetization in intuitionistic PRA}\label{sec:ipra}

\subsection{Primitive-recursive content}

For a fixed coefficient representation, the following functions are primitive
recursive:
\begin{enumerate}[label=(\roman*)]
\item decoding and encoding finite coefficient vectors;
\item addition, Cauchy product, powers, formal derivative, and composition;
\item the fixed-point iteration of Theorem~\ref{thm:fixed};
\item finite Laurent multiplication, residue, and the monomial checks in
Lemma~\ref{lem:res-subst};
\item bounded enumeration of Bell multiplicity vectors and plane rooted trees;
\item finite matrix determinants and multivariate coefficient boxes;
\item rational Taylor partial sums and geometric remainder bounds.
\end{enumerate}
Each item is built from composition and primitive recursion over a bounded list.
The factorial and binomial coefficients in the applications are also primitive
recursive.  Exponentially large values cause no foundational difficulty: PRA
controls computability, not feasible running time.

The cleanest object statement clears all rational denominators.  For example,
Theorem~\ref{thm:lb-der} is represented as
\[
 n\,\operatorname{Coeff}(H\circ G,n)
 =\operatorname{Coeff}(H'\Phi^n,n-1),
\]
using integer numerator/denominator pairs and cross multiplication.  The
conditions ``the input code is a valid reduced rational vector,'' ``the
constant coefficient is zero,'' and ``all equations below stage \(r\) hold''
are decidable primitive-recursive predicates.

\subsection{Quantifier-free induction}

Let \(G_r\) be the coded \(r\)-th fixed-point iterate.  Define
\[
 \operatorname{Inv}(r,N,\Phi)
 =\operatorname{all}_{i<\min(r,N+1)}
   \bigl(\operatorname{Coeff}(G_r,i)=L_i(\Phi)\bigr),
\]
where \(L_i\) denotes the cleared-denominator Lagrange coefficient.  Because
\(\operatorname{all}\), coefficient extraction, and \(L_i\) are function
symbols for primitive-recursive operations, \(\operatorname{Inv}=1\) is a
quantifier-free formula with parameters.  Its base and successor steps are the
finite coefficient calculations of Sections~\ref{sec:finite}--\ref{sec:main}.
Therefore quantifier-free induction proves the uniform assertion for all
\(r\), and instantiating \(r=N+1\) proves the desired truncation theorem.

The same device handles a bounded universal condition in a determinant,
partition, or tree sum.  One never needs to strengthen the induction scheme to
formulas containing unbounded quantifiers.

\subsection{Why the proof is intuitionistic}

No step uses excluded middle for an undecidable statement.  Case splits are on
natural-number equality, rational order, finite-code validity, or an explicitly
supplied nonzero guard, all decidable.  Existence assertions are witnessed by
algorithms: fixed-point iteration, finite reciprocal recursion, Picard
iteration, or a bounded search.  Uniqueness is proved by coefficient agreement
or a rational contraction estimate.  Analytic equality is equality of regular
names with a displayed precision schedule, not equality obtained by choosing a
limit point from a compact set.

\subsection{Analytic names}

For completeness, a regular Gaussian-rational name \(X\) satisfies
\[
 \norm{X(p)-X(q)}_1\le2^{-p}+2^{-q}.
\]
Equality \(X\sim Y\) means
\[
 \norm{X(p)-Y(q)}_1\le2^{-p}+2^{-q}
\]
for all precision indices.  Addition and multiplication have elementary
precision shifts determined by supplied magnitude bounds.  A reciprocal is
formed only with a supplied positive lower bound on the modulus.  The Picard
schedule in \eqref{eq:picard-error} requests
\(q=p+\lceil\log_2\rho\rceil+2\) iterations to obtain precision \(2^{-p}\);
the logarithm is implemented by a bounded search for a power of two.  Thus the
inverse name and its modulus are primitive recursive.

\section{Lean formalization plan}\label{sec:lean}

\subsection{Repository interfaces already available}

The inspected ProveIt snapshot is recorded in \code{SOURCES.md}.  Its
\path{ConciseIPRA/Lean/ConciseIPRA/PR.lean} defines:
\begin{itemize}
\item \code{PR k}, primitive-recursive function symbols of arity \(k\);
\item \code{PTm n}, terms over those symbols;
\item \code{PRAAx}, the defining equations, successor axioms, and
quantifier-free induction;
\item \code{PRA.precZero}, \code{PRA.precSucc}, and \code{PRA.qfInd}.
\end{itemize}
The theorem \code{pra\_theorem\_iff} in
\code{ConciseIPRA.BiInterpretation} transports a closed theorem of the ordinary
\(\ipra\) presentation to the concise scalar theory \(T_{\rm eq}\).

The repository pins Mathlib v4.32.0.  In that version the useful abstract APIs
include \code{PowerSeries.HasSubst}, \code{subst}, \code{coeff\_subst},
\code{substInv}, \code{derivative}, \code{coeff\_derivative},
\code{derivative\_pow}, and \code{derivative\_subst}.  The archive's
\path{lean/LagrangeInversionBlueprint.lean} imports these modules and states the
intended architecture.

\subsection{Recommended module decomposition}

\begin{longtable}{>{\ttfamily}p{0.28\textwidth}p{0.64\textwidth}}
\toprule
Module & Main content\\
\midrule
TruncatedSeries.lean & Finite vectors, bounded Cauchy product, powers,
derivative, composition, coefficient lemmas, and coefficient-ring encodings.\\
FixedPoint.lean & One-degree improvement, stabilization after \(N+1\) steps,
and uniqueness.\\
FormalResidue.lean & Finite Laurent vectors, residue of a derivative, and
change of variables.\\
LagrangeBuermann.lean & Theorems~\ref{thm:lb-alt} and \ref{thm:lb-der}, plus
specializations to \(G\) and \(G^k\).\\
Inverse.lean & Translation to a general base point and connection with
Mathlib's \code{substInv}.\\
Bell.lean & Partial Bell polynomials and Theorem~\ref{thm:bell}.\\
Associahedron.lean & Plane rooted trees, signed tree expansion, and the
correctly normalized face formula.\\
Good.lean & Multivariate truncation, Laurent residue, and determinant formula.\\
Puiseux.lean & Finite binomial roots and ramified branches.\\
Applications.lean & Algebraic equation, Lambert \(W\), shifted series, Catalan
numbers, and exact rational certificates.\\
IPRAEncoding.lean & Natural-number codes and \code{PR} symbols for every
finite algorithm.\\
IPRACorrectness.lean & Quantifier-free invariants and a term of type
\code{Prf PRAAx [] phi}.\\
Transfer.lean & Application of \code{pra\_theorem\_iff}.\\
AnalyticLift.lean & Regular names, contraction bounds, and the explicit inverse
radius; isolated from the algebraic core.\\
\bottomrule
\end{longtable}

\subsection{First abstract targets}

The denominator-free theorem should be formalized first.  In schematic Lean,
with \code{g = X * phi.subst g}, its target is
\begin{lstlisting}[language={}] 
theorem lagrange_burmann_alternative (n : Nat) :
  coeff (n+1) (H.subst g) =
    coeff (n+1)
      (H * phi^n * (phi - X * (d/dX R phi))) := by
  ...
\end{lstlisting}
The exact derivative notation in Mathlib is \code{d\textfractionsolidus dX R}.
The cleared derivative form is
\begin{lstlisting}[language={}]
theorem lagrange_burmann_integral (n : Nat) :
  (n+1 : R) * coeff (n+1) (H.subst g) =
    coeff n ((d/dX R H) * phi^(n+1)) := by
  ...
\end{lstlisting}
The proof should use \code{coeff\_derivative}, \code{derivative\_pow}, and the
Leibniz rule after the alternative theorem.  This order avoids forcing a field
assumption into the core result.

Mathlib's substitution implementation is noncomputable because it represents a
completed power series.  The \(\ipra\) proof must not simply quote that object.
Instead, prove the same identity for \code{TruncSeries}; then establish that
Mathlib coefficient extraction agrees with the finite implementation.  This
separates convenient metatheoretic algebra from the executable object theory.

\subsection{Connection to \code{substInv}}

For a series \(F\) with zero constant coefficient and invertible linear
coefficient, Mathlib already constructs \code{substInv F} and proves the
composition identities.  Define \(\Phi=X/F\) using the unit
\(F/X\), prove
\[
 \operatorname{substInv}(F)=X\,\Phi(\operatorname{substInv}(F)),
\]
and specialize the abstract theorem.  The Catalan example can then be connected
to existing Mathlib Catalan power-series lemmas, providing a useful integration
test.

\subsection{Object-theory target}

After coding the algorithms, the central deliverable is not another meta-level
Lean equality but a proof object
\begin{lstlisting}[language={}]
Prf PRAAx [] lagrangeCorrectnessFormula
\end{lstlisting}
where \code{lagrangeCorrectnessFormula : Fm PTm 0} is the closed arithmetic
correctness statement.  Transport is then exactly
\begin{lstlisting}[language={}]
(pra_theorem_iff lagrangeCorrectnessFormula).mp proof
\end{lstlisting}
as shown in the blueprint.  Replacing its two \code{opaque} placeholders by
this derivation is the decisive completion criterion.

\section{Validation, limitations, and reproducibility}\label{sec:validation}

The exact-rational program \code{code/verify.py} independently checks:
\begin{itemize}
\item both univariate Lagrange--B\"urmann forms through degree nine for
nontrivial rational input series;
\item the Bell-polynomial formula through degree eight;
\item the equation \(x^p-x+z=0\) for \(2\le p\le6\) through degree 24;
\item Lambert \(W\) through degree 12 and the shifted series through degree six;
\item Catalan coefficients through degree 15;
\item a bivariate instance of Good's formula on a \(4\times4\) coefficient box;
\item the rational interval certificate \eqref{eq:omega-cert};
\item the degree-two associahedron normalization correction.
\end{itemize}
The recorded output ends with \code{ALL EXACT CHECKS PASSED}.  These tests are
not substitutes for proof, but they exercise the finite formulae and catch
index, factorial, sign, and Jacobian transposition errors.

The Lean file is a formalization blueprint, not a completed kernel-checked
proof.  In particular, the exact analytic continuation argument for the shifted
Lambert radius is the heaviest remaining Lean component.  The algebraic
coefficient theorem, fixed-point iteration, Bell formula, and finite
applications require only finite sums and should be completed first.

\section{Further research questions}\label{sec:further}

\begin{enumerate}
\item Can the finite residue substitution lemma be packaged as a reusable
\(\ipra\) library theorem for other coefficient-extraction arguments, such as
Good inversion, diagonal extraction, and the residue theorem for rational
functions?
\item What is the smallest quantifier-free induction invariant that proves the
full Lagrange--B\"urmann identity after denominator clearing?  A direct
coefficient invariant may be substantially smaller than an explicit Laurent
calculus.
\item Can the associahedron proof be formalized through a cancellation-free
antipode theorem for a finite Hopf algebra of faces, and can that proof be made
shorter than the rooted-tree expansion in the \code{ConciseIPRA} checker
metric?
\item For the shifted Lambert series, can one derive a purely coefficientwise
proof of the exact radius \(\sqrt{4+\pi^2}\), avoiding effective covering-space
machinery?
\item How far can the ramified construction be pushed toward a fully effective
Newton--Puiseux theorem in \(\ipra\), with explicit denominator and coefficient
bounds?
\item Can the same finite formal-residue core prove the multivariate
Lagrange--Good formula uniformly in a coded dimension without a prohibitive
proof-size increase?
\end{enumerate}

\section{Conclusion}

The essential content of Lagrange inversion is finite.  A coefficient of degree
\(n\) is determined by at most degree \(n\) data, the implicit solution is
created by \(n+1\) fixed-point iterations, and the residue proof reduces to a
finite Laurent-polynomial identity.  This makes the theorem especially well
suited to intuitionistic primitive recursive arithmetic.  Analytic existence
requires additional quantitative data, but an explicit contraction supplies a
primitive-recursive inverse name and a positive radius.  The page's principal
examples and extensions then follow by finite algebra, with exact convergence
arguments supplied where the page asserts radii.

\begin{thebibliography}{99}

\bibitem{WikipediaLagrange}
Wikipedia contributors,
``Lagrange inversion theorem,'' revision 1348379317, April 12, 2026.
\url{https://en.wikipedia.org/w/index.php?title=Lagrange_inversion_theorem&oldid=1348379317}.

\bibitem{Gessel}
Ira M. Gessel,
``Lagrange inversion,''
\emph{Journal of Combinatorial Theory, Series A} 144 (2016), 212--249.
\url{https://arxiv.org/abs/1609.05988}.

\bibitem{SuryaWarnke}
Erlang Surya and Lutz Warnke,
``Lagrange Inversion Formula by Induction,''
\emph{American Mathematical Monthly} 130 (2023), 944--948.
\url{https://arxiv.org/abs/2305.17576}.

\bibitem{AguiarArdila}
Marcelo Aguiar and Federico Ardila,
\emph{Hopf Monoids and Generalized Permutahedra},
Memoirs of the American Mathematical Society 1437.
\url{https://arxiv.org/abs/1709.07504}.

\bibitem{CorlessJeffreyKnuth}
Robert M. Corless, David J. Jeffrey, and Donald E. Knuth,
``A sequence of series for the Lambert W function,''
ISSAC 1997, 197--204.

\bibitem{proveit}
Vladimir Reshetnikov,
\emph{ProveIt}, repository snapshot recorded in \code{SOURCES.md}.
\url{https://github.com/VladimirReshetnikov/ProveIt}.

\bibitem{proveitTaylor}
Vladimir Reshetnikov,
``Taylor's theorem and Cauchy's integral formula in intuitionistic primitive
recursive arithmetic,'' ProveIt proof packet,
\code{Logic/PrimitiveRecursiveArithmetic/docs/proofs/analysis/taylor\_cauchy\_ipra/}.

\bibitem{mathlibSubst}
Mathlib contributors,
\code{Mathlib.RingTheory.PowerSeries.Substitution}, v4.32.0.
\url{https://github.com/leanprover-community/mathlib4/blob/v4.32.0/Mathlib/RingTheory/PowerSeries/Substitution.lean}.

\bibitem{mathlibDerivative}
Mathlib contributors,
\code{Mathlib.RingTheory.PowerSeries.Derivative}, v4.32.0.
\url{https://github.com/leanprover-community/mathlib4/blob/v4.32.0/Mathlib/RingTheory/PowerSeries/Derivative.lean}.

\end{thebibliography}

\end{document}
